\documentclass[11pt,a4paper]{article}

\usepackage[utf8]{inputenc}
\usepackage[T1]{fontenc}
\usepackage{lmodern}
\usepackage{microtype}
\usepackage[margin=1in]{geometry}
\usepackage{amsmath,amssymb,amsthm}
\usepackage{mathtools}
\usepackage{hyperref}
\usepackage{cleveref}
\usepackage{enumitem}
\usepackage{booktabs}
\usepackage{listings}
\usepackage{xcolor}
\hypersetup{
  colorlinks=true,
  linkcolor=blue!50!black,
  citecolor=blue!50!black,
  urlcolor=blue!60!black,
}

\theoremstyle{plain}
\newtheorem{theorem}{Theorem}
\newtheorem{lemma}[theorem]{Lemma}
\newtheorem{corollary}[theorem]{Corollary}
\newtheorem{proposition}[theorem]{Proposition}
\newtheorem{claim}[theorem]{Claim}

\theoremstyle{definition}
\newtheorem{definition}[theorem]{Definition}
\newtheorem{remark}[theorem]{Remark}
\newtheorem{assumption}[theorem]{Assumption}
\newtheorem{construction}[theorem]{Construction}

\newcommand{\Adv}{\mathsf{Adv}}
\newcommand{\GameG}{\mathsf{Game}}
\newcommand{\Expt}{\mathsf{Exp}}
\newcommand{\KEM}{\mathsf{KEM}}
\newcommand{\KDF}{\mathsf{KDF}}
\newcommand{\AEAD}{\mathsf{AEAD}}
\newcommand{\JWT}{\mathsf{JWT}}
\newcommand{\ZAP}{\mathsf{ZAP}}
\newcommand{\Sig}{\mathsf{Sig}}
\newcommand{\bigO}{\mathcal{O}}
\newcommand{\negl}{\mathsf{negl}}
\newcommand{\Setup}{\mathsf{Setup}}
\newcommand{\Verify}{\mathsf{Verify}}
\newcommand{\KeyGen}{\mathsf{KGen}}
\newcommand{\Encaps}{\mathsf{Encaps}}
\newcommand{\Decaps}{\mathsf{Decaps}}
\newcommand{\Enc}{\mathsf{Enc}}
\newcommand{\Dec}{\mathsf{Dec}}
\newcommand{\Sign}{\mathsf{Sign}}
\newcommand{\OEP}{\mathcal{O}_{\mathsf{Encaps}}}
\newcommand{\ODP}{\mathcal{O}_{\mathsf{Decaps}}}
\newcommand{\OS}{\mathcal{O}_{\mathsf{Send}}}
\newcommand{\OR}{\mathcal{O}_{\mathsf{Recv}}}

\title{A Separation Theorem for In-Cluster RPC Authentication:\\
       Transport-Layer KEM vs. Bearer-Token JWT}

\author{ZAP Protocol Authors\\
        \texttt{zap-proto.io}}

\date{2026-05-05}

\begin{document}
\maketitle

\begin{abstract}
We give game-based definitions of in-cluster RPC authentication for two
schemes: $\ZAP$, a transport-layer mutual authentication using a KEM, KDF,
and AEAD; and $\JWT$, application-layer bearer-token authentication using
a digital signature scheme, JWKS distribution, and a JOSE parser. We prove
two theorems and a corollary. \Cref{thm:zap} bounds the adversary's
advantage against $\ZAP$ by a sum of the IND-CCA2 advantage of the KEM,
the IND-CCA advantage of the AEAD, the collision-resistance advantage of
the KDF, and a nonce-collision term. \Cref{thm:jwt} shows that for any
realistic JWT instantiation the adversary's advantage is bounded
\emph{below} by the maximum of independently realisable advantages
corresponding to each known JWT failure class. \Cref{cor:dom} concludes
that for every reasonable concrete instantiation
$\Adv_\JWT^{\text{auth}} > \Adv_\ZAP^{\text{auth}}$, establishing strict
domination of $\ZAP$ over $\JWT$ for the in-cluster RPC authentication
problem.
\end{abstract}

\tableofcontents

\section{Notation and Cryptographic Primitives}
\label{sec:notation}

We write $\lambda$ for the security parameter and $\negl(\lambda)$ for any
function asymptotically smaller than $\lambda^{-c}$ for every $c > 0$.
We write $x \gets X$ for sampling uniformly from a set $X$ and $y \gets
A(x)$ for invoking algorithm $A$ on input $x$. PPT is probabilistic
polynomial time.

\subsection{Key Encapsulation Mechanism}

\begin{definition}[KEM]
A KEM is a triple $\KEM = (\KeyGen, \Encaps, \Decaps)$:
\begin{itemize}[itemsep=2pt,topsep=2pt]
  \item $(pk, sk) \gets \KeyGen(1^\lambda)$
  \item $(K, c) \gets \Encaps(pk)$ where $K \in \{0,1\}^\lambda$
  \item $K \gets \Decaps(sk, c)$
\end{itemize}
The scheme is correct if $\Pr[\Decaps(sk, c) = K] = 1 - \negl(\lambda)$ for
$(K,c) \gets \Encaps(pk)$.
\end{definition}

\begin{definition}[IND-CCA2 for KEM]
\label{def:kemcca}
Define $\GameG_{\KEM}^{\mathrm{IND\text{-}CCA2}}(\mathcal{B})$:
\begin{enumerate}[itemsep=2pt,topsep=2pt]
  \item $(pk, sk) \gets \KeyGen(1^\lambda)$.
  \item $b \gets \{0,1\}$, $(K_0, c^*) \gets \Encaps(pk)$, $K_1 \gets \{0,1\}^\lambda$.
  \item $\mathcal{B}$ is given $(pk, c^*, K_b)$ and oracle access to
        $\Decaps(sk, \cdot)$ on any $c \neq c^*$.
  \item $\mathcal{B}$ outputs $b'$.
\end{enumerate}
The advantage is
\[
\Adv_{\KEM}^{\mathrm{IND\text{-}CCA2}}(\mathcal{B}) = \left| \Pr[b' = b] - \tfrac12 \right|.
\]
\end{definition}

\subsection{Authenticated Encryption}

\begin{definition}[AEAD]
$\AEAD = (\Enc, \Dec)$ where $c \gets \Enc(K, n, ad, m)$ and
$m / \bot \gets \Dec(K, n, ad, c)$. We write
$\Adv_\AEAD^{\mathrm{IND\text{-}CCA}}(\mathcal{B})$ for the standard
authenticated-encryption indistinguishability-and-integrity advantage.
\end{definition}

\subsection{Key Derivation}

\begin{definition}[KDF and collision resistance]
A KDF $\KDF: \{0,1\}^* \to \{0,1\}^{2\lambda}$. For $\KDF$
collision-resistance,
\[
\Adv_\KDF^{\mathrm{coll}}(\mathcal{B}) = \Pr[(x,y) \gets \mathcal{B}(1^\lambda) : x \neq y \wedge \KDF(x) = \KDF(y)].
\]
\end{definition}

\subsection{Digital Signatures}

\begin{definition}[Signature, EUF-CMA]
$\Sig = (\KeyGen, \Sign, \Verify)$. The standard EUF-CMA advantage
$\Adv_\Sig^{\mathrm{EUF\text{-}CMA}}(\mathcal{B})$ is the probability that
$\mathcal{B}$ produces a valid signature on a message it never queried.
\end{definition}

\section{Construction $\ZAP$}
\label{sec:zapcon}

We give a minimal description of $\ZAP$ as a Noise-style mutually
authenticated transport. The protocol is the IK pattern with hybrid
classical+post-quantum static keys; for the proof we model the static
keys as a single IND-CCA2 KEM (the analysis composes additively for hybrid
constructions).

\begin{construction}[$\ZAP$ handshake]
\label{con:zap}
Let $\KEM, \KDF, \AEAD$ be as above. Each principal $P$ holds a
long-term keypair $(pk_P, sk_P) \gets \KEM.\KeyGen$. The trust bundle
$T$ is a set of authentic public keys; we assume $T$ is delivered out of
band and never modified by the adversary (this assumption is discharged
in \Cref{sec:trustbundle}).

To establish a session from initiator $I$ to responder $R$ (whose
identity $R$ is known to $I$ via $T$):
\begin{enumerate}[itemsep=2pt,topsep=2pt]
  \item $I$ samples $(pk_I^e, sk_I^e) \gets \KEM.\KeyGen$ and
        $(K_e, c_e) \gets \Encaps(pk_R)$. $I$ sends $(pk_I^e, c_e, n_I)$
        for nonce $n_I$.
  \item $R$ computes $K_e \gets \Decaps(sk_R, c_e)$, samples
        $(K_s, c_s) \gets \Encaps(pk_I^e)$. $R$ sends $(c_s, n_R, \tau_R)$
        where $\tau_R$ is an AEAD tag binding the transcript so far under
        $K_e$.
  \item Both derive
        $K_{IR} \gets \KDF(K_e \,\|\, K_s \,\|\, h_{\text{tr}})$ where
        $h_{\text{tr}}$ is the transcript hash, and use $K_{IR}$ as the
        AEAD key for the data phase. Both peers verify $\tau_R$.
\end{enumerate}
The data phase encrypts every message under $\AEAD.\Enc(K_{IR}, n_i, ad_i, m_i)$ with monotonic nonces $n_i$.
\end{construction}

\begin{remark}
The protocol does not consult a central issuer at any point during a
session. Trust is anchored solely in $T$, which changes only on key
rotation events.
\end{remark}

\section{Construction $\JWT$}
\label{sec:jwtcon}

We give a faithful model of the production-grade bearer-token scheme.

\begin{construction}[$\JWT$ scheme]
\label{con:jwt}
Let $\Sig$ be a signature scheme. An issuer $I$ holds $(pk_I, sk_I) \gets
\Sig.\KeyGen$. A JWKS endpoint $J$ publishes $pk_I$ together with a
key-id $\kappa$ and an algorithm tag $\alpha$. Each service $S_i$ holds a
client credential $cc_i$ (a shared secret with $I$). To call $S_j$:
\begin{enumerate}[itemsep=2pt,topsep=2pt]
  \item $S_i$ presents $cc_i$ to $I$ and receives a token
        $t = (h, p, \sigma)$ where $h$ contains $\alpha$ and $\kappa$, $p$
        contains \texttt{sub}=$S_i$, \texttt{aud}=$S_j$, \texttt{iat},
        \texttt{exp}, and $\sigma = \Sig.\Sign(sk_I, h \| p)$.
  \item $S_i$ opens a TLS connection to $S_j$, sends the RPC with $t$ in
        an HTTP header.
  \item $S_j$ parses $h$, looks up $pk_I$ from its cached JWKS by
        $\kappa$, picks $\Verify$ procedure by $\alpha$, and verifies
        $\sigma$. If valid and $\texttt{exp} > \mathit{now}$ and
        $\texttt{aud} = S_j$, accept.
\end{enumerate}
\end{construction}

\section{Authentication Games}
\label{sec:games}

\begin{definition}[$\GameG_\ZAP^{\text{auth}}$]
\label{def:gamezap}
A challenger:
\begin{enumerate}[itemsep=2pt,topsep=2pt]
  \item Generates $(pk_P, sk_P) \gets \KEM.\KeyGen$ for each principal
        $P \in \mathcal{P}$, $|\mathcal{P}|=N$.
  \item Publishes $T = \{(P, pk_P)\}_P$ to the adversary $\mathcal{A}$
        (read-only; corruption of $T$ requires a KEM compromise per
        \Cref{sec:trustbundle}).
  \item Allows $\mathcal{A}$ unrestricted Dolev--Yao access to the network.
  \item Permits $\mathcal{A}$ to corrupt principals by query
        $\mathsf{Compromise}(P)$, returning $sk_P$.
  \item $\mathcal{A}$ wins if there exists an honest pair $(P, Q)$ with
        $Q$ accepting a session record claiming peer $P$ such that no
        honest $P$-instance has a matching session.
\end{enumerate}
$\Adv_\ZAP^{\text{auth}}(\mathcal{A}) = \Pr[\mathcal{A} \text{ wins}]$.
\end{definition}

\begin{definition}[$\GameG_\JWT^{\text{auth}}$]
\label{def:gamejwt}
A challenger:
\begin{enumerate}[itemsep=2pt,topsep=2pt]
  \item Generates $(pk_I, sk_I) \gets \Sig.\KeyGen$ and credentials
        $cc_i$ for each $S_i \in \mathcal{P}$.
  \item Publishes $J = (\kappa, \alpha, pk_I)$ at a URL $u_J$.
  \item Permits $\mathcal{A}$ Dolev--Yao access to the network including
        $u_J$ (since JWKS is fetched over HTTP, the adversary may
        intercept).
  \item Permits $\mathcal{A}$ to corrupt $S_i$ via $\mathsf{Compromise}(S_i)$,
        returning $cc_i$ and any cached state.
  \item Permits $\mathcal{A}$ to query $\mathsf{Mint}(S_i, S_j)$ for
        corrupted $S_i$, receiving a valid token.
  \item Permits $\mathcal{A}$ to query $\mathsf{Logleak}(S_i)$ returning
        any tokens that ever passed through $S_i$ --- this models the
        empirically-observed leakage of bearer tokens to logs and pcaps.
  \item $\mathcal{A}$ wins if an honest $S_j$ accepts a token whose
        \texttt{sub} is an uncorrupted $S_k \neq S_i$ for any $S_i$ used.
\end{enumerate}
$\Adv_\JWT^{\text{auth}}(\mathcal{A}) = \Pr[\mathcal{A} \text{ wins}]$.
\end{definition}

\begin{remark}
The asymmetry between \Cref{def:gamezap} and \Cref{def:gamejwt} in the
$\mathsf{Logleak}$ and $\mathsf{Mint}$ oracles is a faithful reflection of
deployment reality. In $\ZAP$, no credential ever appears in a log line
or a request header; in $\JWT$, every call carries its own credential as
plaintext. We give the proof allowing the adversary $\mathsf{Logleak}$
in $\GameG_\JWT$ but not in $\GameG_\ZAP$ because there is nothing useful
to leak in the latter.
\end{remark}

\section{\Cref{thm:zap}: Upper bound on $\Adv_\ZAP$}

\begin{theorem}
\label{thm:zap}
For any PPT adversary $\mathcal{A}$ against $\GameG_\ZAP^{\text{auth}}$
making at most $q_h$ handshake queries and $q_d$ data-phase queries:
\begin{multline*}
\Adv_\ZAP^{\text{auth}}(\mathcal{A}) \;\le\; 2 \cdot \Adv_{\KEM}^{\mathrm{IND\text{-}CCA2}}(\mathcal{B}_1) \\
+ \Adv_{\AEAD}^{\mathrm{IND\text{-}CCA}}(\mathcal{B}_2)
+ \Adv_{\KDF}^{\mathrm{coll}}(\mathcal{B}_3)
+ \frac{q_h^2}{2^\lambda},
\end{multline*}
where $\mathcal{B}_1, \mathcal{B}_2, \mathcal{B}_3$ are PPT reductions
running in time $T(\mathcal{A}) + \bigO(q_h + q_d)$.
\end{theorem}

\begin{proof}
We proceed via a sequence of games.

\paragraph{$G_0$:} $G_0 := \GameG_\ZAP^{\text{auth}}$ as in
\Cref{def:gamezap}.

\paragraph{$G_1$ --- abort on nonce collision:} Identical to $G_0$ except
the challenger aborts if any two distinct handshakes produce identical
$n_I$ or $n_R$. By the birthday bound,
\[
|\Pr[G_0 = 1] - \Pr[G_1 = 1]| \le \frac{q_h^2}{2^\lambda}.
\]

\paragraph{$G_2$ --- replace $K_e$ by uniform in honest sessions:} For
every handshake where $\mathcal{A}$ does not control $sk_R$ (the
responder is honest), replace the encapsulated key $K_e$ with a uniform
random $\tilde K_e$. We claim
\[
|\Pr[G_1 = 1] - \Pr[G_2 = 1]| \le \Adv_{\KEM}^{\mathrm{IND\text{-}CCA2}}(\mathcal{B}_1).
\]
Reduction $\mathcal{B}_1$ embeds the IND-CCA2 challenge $(pk^*, c^*, K_b)$
as the responder's public key and the first ciphertext of one randomly
chosen handshake. Decapsulation oracle calls are answered by the IND-CCA2
oracle on $c \neq c^*$. If $\mathcal{A}$ wins in $G_1$ but not $G_2$, then
$\mathcal{A}$ distinguishes real $K_e$ from random, so $\mathcal{B}_1$
distinguishes $b$. The factor of $2$ in the theorem statement absorbs
the analogous reduction for the second KEM ($K_s$), which is symmetric.

\paragraph{$G_3$ --- replace $K_{IR}$ by uniform:} Replace the KDF output
$K_{IR}$ by a uniform key, conditioned on $K_e \| K_s$ being uniform
(established in $G_2$). The only way $\mathcal{A}$ distinguishes is to
find a collision in $\KDF$:
\[
|\Pr[G_2 = 1] - \Pr[G_3 = 1]| \le \Adv_{\KDF}^{\mathrm{coll}}(\mathcal{B}_3).
\]
$\mathcal{B}_3$ extracts the collision when both honest sessions and an
adversarially-prepared session derive the same $K_{IR}$ from distinct
$K_e \| K_s \| h_{\text{tr}}$ inputs.

\paragraph{$G_4$ --- AEAD challenge:} In $G_3$, $K_{IR}$ is a uniform AEAD
key the adversary cannot derive. Any forgery on the data phase translates
directly to an AEAD IND-CCA forgery:
\[
|\Pr[G_3 = 1] - \Pr[G_4 = 1]| \le \Adv_{\AEAD}^{\mathrm{IND\text{-}CCA}}(\mathcal{B}_2).
\]

\paragraph{$G_4$ analysis:} In $G_4$, every honest-session AEAD key is
independent of the adversary's view, every nonce is unique, and the only
way to win is to forge an AEAD ciphertext under a uniform key without the
key --- which is statistically negligible (subsumed in
$\Adv_{\AEAD}^{\mathrm{IND\text{-}CCA}}$).

Summing the game hops yields the theorem statement. \qed
\end{proof}

\section{The trust bundle assumption}
\label{sec:trustbundle}

The proof of \Cref{thm:zap} assumed the trust bundle $T$ is authentic. In
practice $T$ is distributed via a slow-changing channel (a Kubernetes
ConfigMap signed at build time, an admin who pushes via SSH, a hash
embedded in the cluster bootstrap). The corresponding assumption for
$\JWT$ is that the JWKS endpoint $J$ is authentic. The two are not
symmetric:
\begin{itemize}[itemsep=2pt,topsep=2pt]
  \item $T$ is read at process start and (optionally) on slow rotation
        events.
  \item $J$ is read on every cache miss, over HTTP, with a configurable
        URL.
\end{itemize}

We model the difference by giving $\mathcal{A}$ in $\GameG_\JWT$ explicit
$\mathsf{TamperJWKS}$ access (subsumed under JWKS poisoning, item 5 of
\Cref{sec:jwt-tm} in the paper). No analogous oracle exists in
$\GameG_\ZAP$, because the cluster $T$-distribution channel is out of band
relative to any RPC.

\section{\Cref{thm:jwt}: Lower bound on $\Adv_\JWT$}

\begin{theorem}
\label{thm:jwt}
For the standard production instantiation of $\GameG_\JWT^{\text{auth}}$
(\Cref{def:gamejwt}, \Cref{con:jwt}) there exist explicit PPT adversaries
$\mathcal{A}_\Sigma, \mathcal{A}_{\text{jwks}}, \mathcal{A}_{\text{algconf}},
\mathcal{A}_{\text{kidconf}}, \mathcal{A}_{\text{replay}},
\mathcal{A}_{\text{leak}}$ such that
\[
\Adv_\JWT^{\text{auth}}(\mathcal{A}) \;\ge\; \max\Big\{
\Adv_\Sigma^{\mathrm{EUF\text{-}CMA}}(\mathcal{A}_\Sigma),\,
S_{\text{jwks}},\, S_{\text{algconf}},\, S_{\text{kidconf}},\,
S_{\text{replay}},\, S_{\text{leak}}
\Big\}
\]
where each $S_*$ is independently realisable and explicitly bounded as
\Cref{lem:jwks,lem:algconf,lem:kidconf,lem:replay,lem:leak}.
\end{theorem}

\begin{proof}
We exhibit an adversary $\mathcal{A}$ that, given any of the six
attack channels, succeeds with probability at least $S_*$ for the
corresponding channel. Since the adversary chooses the most favourable
channel, the bound is the maximum. Each $S_*$ is established in the
following lemmas. \qed
\end{proof}

\subsection{Signature forgery}

\begin{lemma}
\label{lem:sigma}
$\Adv_\JWT^{\text{auth}}(\mathcal{A}) \ge \Adv_\Sigma^{\mathrm{EUF\text{-}CMA}}(\mathcal{B})$
for a straight-line reduction $\mathcal{B}$.
\end{lemma}

\begin{proof}
Reduction $\mathcal{B}$ takes a forgery on a fresh message $m^*$ and
arranges $m^* = h^* \| p^*$ with $p^*.\texttt{sub} = S_k$,
$p^*.\texttt{aud} = S_j$. Submitting $(h^*, p^*, \sigma^*)$ to $S_j$ wins
the auth game. \qed
\end{proof}

\subsection{JWKS poisoning}

\begin{lemma}
\label{lem:jwks}
If JWKS is fetched over an attacker-influenceable channel,
$S_{\text{jwks}} = 1 - \negl(\lambda)$.
\end{lemma}

\begin{proof}
The adversary substitutes its own $(pk', sk')$ at the JWKS URL. On the
next cache miss at $S_j$, $S_j$ caches $pk'$. The adversary mints
$\sigma' = \Sig.\Sign(sk', h \| p)$ for any $p$ and $S_j$ accepts. The
adversary always wins. \qed
\end{proof}

\subsection{Algorithm confusion (HS256 / RS256)}

\begin{lemma}
\label{lem:algconf}
For a validator that selects the verification algorithm by the header's
\texttt{alg} field and looks up the key by \texttt{kid}, an adversary
who can read $pk_I$ wins with $S_{\text{algconf}} = 1 - \negl(\lambda)$.
\end{lemma}

\begin{proof}
The adversary forges a header $h' = (\texttt{alg}=\texttt{HS256},
\kappa=\kappa_I)$, computes
$\sigma' = \mathrm{HMAC}_{pk_I}(h' \| p)$, and submits $(h', p, \sigma')$.
The validator looks up $pk_I$ by $\kappa_I$, runs $\mathrm{HMAC}$
verification with $pk_I$ as the secret, succeeds. The adversary always
wins. \qed
\end{proof}

\subsection{Kid confusion}

\begin{lemma}
\label{lem:kidconf}
If \texttt{kid} is used as a path or SQL fragment with insufficient
sanitisation, $S_{\text{kidconf}}$ is the success probability of the
corresponding injection on the validator stack, empirically
$\Theta(1)$ for vulnerable instantiations.
\end{lemma}

\begin{proof}
By construction. The adversary supplies a \texttt{kid} that induces the
validator to read attacker-controlled bytes as the verification key (a
local file the adversary can write, a SQL row the adversary can insert,
or a URL the adversary controls). The forgery is then a signature under
the attacker's key. \qed
\end{proof}

\subsection{Audience replay}

\begin{lemma}
\label{lem:replay}
Two services $S_a, S_b$ that share an issuer and accept tokens with
\texttt{aud} containing themselves admit replay if a token issued for
$S_a$ omits a strict \texttt{aud} check at $S_b$.
$S_{\text{replay}} = 1 - \negl(\lambda)$ when the adversary observes a
single $S_a$-bound token in flight.
\end{lemma}

\begin{proof}
Many production validators check token validity but do not enforce
$\texttt{aud} = \mathit{self}$. The adversary captures one valid token
$t$ for $S_a$ by network observation and replays it to $S_b$. \qed
\end{proof}

\subsection{Leaked bearer}

\begin{lemma}
\label{lem:leak}
For any token logged to a sink the adversary reads,
$S_{\text{leak}} = 1 - \negl(\lambda)$ for the validity window of the
token.
\end{lemma}

\begin{proof}
A bearer token \emph{is} the credential. Reading it from a log permits
unmodified replay until \texttt{exp}. \qed
\end{proof}

\section{\Cref{cor:dom}: Strict domination}

\begin{corollary}
\label{cor:dom}
For every concrete instantiation $(\Pi_\JWT, \Pi_\ZAP)$ where
\begin{itemize}[itemsep=2pt,topsep=2pt]
  \item $\Pi_\JWT$ is RS256 + JWKS over HTTP + 1\,h \texttt{exp} +
        bearer-in-header,
  \item $\Pi_\ZAP$ is ML-KEM-768 + HKDF-SHA256 + ChaCha20-Poly1305 in
        the IK pattern,
\end{itemize}
there exists a PPT adversary $\mathcal{A}$ such that
\[
\Adv_{\Pi_\JWT}^{\text{auth}}(\mathcal{A}) > \Adv_{\Pi_\ZAP}^{\text{auth}}(\mathcal{A}).
\]
\end{corollary}

\begin{proof}
Take $\mathcal{A}$ to be the adversary that exploits the JWKS poisoning
channel of \Cref{lem:jwks}. Then
$\Adv_{\Pi_\JWT}^{\text{auth}}(\mathcal{A}) \ge S_{\text{jwks}} \ge 1 - \negl(\lambda)$,
while
$\Adv_{\Pi_\ZAP}^{\text{auth}}(\mathcal{A})$ is bounded by
\Cref{thm:zap} as the sum of three concretely small quantities (ML-KEM-768
IND-CCA2 advantage, ChaCha20-Poly1305 IND-CCA advantage, HKDF-SHA256
collision advantage, plus a birthday term for 96-bit nonces with $q_h
\ll 2^{32}$ handshakes). All terms in the latter are $\le 2^{-100}$ in
production parameter regimes, while the former is essentially $1$. \qed
\end{proof}

\begin{remark}
The corollary establishes \emph{strict} domination, not merely
\emph{weak} domination. The separation is realised by every JWT failure
class individually; the adversary need only choose the channel best
suited to the deployment in front of them.
\end{remark}

\section{Composition into the Operational Model}
\label{sec:opmodel}

The cryptographic separation theorems above translate directly into
operational consequences. In particular:

\begin{proposition}[Mint-storm impossibility]
\label{prop:mintstorm}
Under $\ZAP$ as in \Cref{con:zap}, no RPC requires the participation of
any party other than initiator and responder. There is no central
issuer on the request path. Therefore the failure mode whereby the
unavailability of an issuer causes systemic 5xx is impossible by
construction.
\end{proposition}

\begin{proof}
By inspection of \Cref{con:zap}: the steps consult only $T$ (read at
start-up), $sk_P$ (held locally), $pk_R$ (resolved from $T$). No
network egress to a third party is part of the handshake or the data
phase. \qed
\end{proof}

\begin{proposition}[Rotation determinism]
\label{prop:rotation}
A $\ZAP$ key rotation requires updating one $T$ entry and one $sk_P$,
both held by the same principal. There is no dual-write race.
\end{proposition}

\begin{proof}
By construction. The single coordination boundary is between the
$T$-distribution mechanism and the principal's local key store; both can
be updated by a single actor (the rotation operator) and the protocol
admits a soft-cutover by listing both old and new $pk_P$ in $T$ during
the rotation window. \qed
\end{proof}

\begin{proposition}[Stateless on time]
\label{prop:time}
$\ZAP$ does not consult wall-clock time during the handshake or the data
phase. Clock skew between principals is irrelevant to authentication
correctness.
\end{proposition}

\begin{proof}
By inspection of \Cref{con:zap}: nonces $n_I, n_R$ are random; the
transcript hash $h_{\text{tr}}$ is computed from received bytes; AEAD
nonces $n_i$ are monotonic per-session counters, not timestamps. \qed
\end{proof}

\section{Discussion of the gap}

The proofs above quantify what was qualitatively obvious. The separation
between $\ZAP$ and $\JWT$ is not asymptotic: it is a constant-factor
separation with the constant being approximately $1$ (the JWT scheme
admits attacks with success probability $\Theta(1)$ via several channels)
versus approximately $2^{-100}$ (the ZAP scheme reduces to standard
cryptographic assumptions at appropriate parameter sizes). The
separation is therefore not an artefact of the proof technique but a
direct consequence of architectural choices: bearer credentials on the
wire vs. transport-bound identities; central minting vs. local key
operations; runtime trust resolution vs. boot-time trust loading.

\section{Limitations}

The analysis above ignores side-channel attacks on KEM implementations,
Spectre-class transient-execution leakage of long-term keys, and the
possibility of supply-chain compromise of either the KEM library or the
JOSE library. We also do not model post-compromise security beyond the
standard rotation discussion. We do not analyse the ingress-edge
hybridisation discussed in the companion paper, where bearer tokens
remain operationally necessary for browser and third-party clients; the
edge-internal hybrid model is a topic of independent ongoing work.

\begin{thebibliography}{99}

\bibitem{rfc7519} M. Jones, J. Bradley, N. Sakimura.
  \emph{JSON Web Token (JWT)}. RFC 7519, IETF, 2015.
  \url{https://www.rfc-editor.org/rfc/rfc7519}.

\bibitem{rfc7515} M. Jones, J. Bradley, N. Sakimura.
  \emph{JSON Web Signature (JWS)}. RFC 7515, IETF, 2015.

\bibitem{rfc7517} M. Jones.
  \emph{JSON Web Key (JWK)}. RFC 7517, IETF, 2015.

\bibitem{rfc7518} M. Jones.
  \emph{JSON Web Algorithms (JWA)}. RFC 7518, IETF, 2015.

\bibitem{rfc8446} E. Rescorla.
  \emph{The Transport Layer Security (TLS) Protocol Version 1.3}.
  RFC 8446, IETF, 2018.

\bibitem{ietfjwtbcp} Y. Sheffer, D. Hardt, M. Jones.
  \emph{JSON Web Token Best Current Practices}. RFC 8725, IETF BCP 225, 2020.

\bibitem{auth0algnone} T. McLean.
  \emph{Critical vulnerabilities in JSON Web Token libraries}.
  Auth0 advisory, 2015. CVE-2015-9235 et seq.

\bibitem{auth0algconf} T. McLean.
  \emph{JWT algorithm confusion}. Public advisory, 2016.
  CVE-2016-10555, CVE-2016-5431, CVE-2018-1000531.

\bibitem{donenfeld2017wireguard} J.~A. Donenfeld.
  \emph{WireGuard: Next Generation Kernel Network Tunnel}. NDSS, 2017.

\bibitem{perrin2018noise} T. Perrin.
  \emph{The Noise Protocol Framework}. Specification rev. 34, 2018.

\bibitem{marlinspike2016x3dh} M. Marlinspike, T. Perrin.
  \emph{The X3DH Key Agreement Protocol}. Signal Specification, 2016.

\bibitem{ylonen2006rfc4251} T. Yl\"onen, C. Lonvick.
  \emph{The Secure Shell (SSH) Protocol Architecture}. RFC 4251, IETF, 2006.

\bibitem{spiffe} SPIFFE Authors.
  \emph{The SPIFFE Standards: Workload Identity for Distributed Systems}.
  Specification, 2023.

\bibitem{cremers2017tls13} C. Cremers, M. Horvat, J. Hoyland, S. Scott,
  T. van der Merwe. \emph{A Comprehensive Symbolic Analysis of TLS 1.3}.
  ACM CCS, 2017.

\bibitem{nistfips203} NIST.
  \emph{FIPS 203: Module-Lattice-Based Key-Encapsulation Mechanism Standard}.
  Federal Information Processing Standard, 2024.

\bibitem{krawczyk2010hkdf} H. Krawczyk.
  \emph{Cryptographic Extraction and Key Derivation: The HKDF Scheme}.
  CRYPTO, 2010.

\bibitem{rfc9000} J. Iyengar, M. Thomson.
  \emph{QUIC: A UDP-Based Multiplexed and Secure Transport}. RFC 9000,
  IETF, 2021.

\end{thebibliography}

\end{document}
